\topic{Introduction}{A moving lid, a circulating fluid and an unknown pressure}
Consider a square container filled with a viscous liquid. Its top wall moves
to the right while the other three walls remain stationary. Friction with
the moving wall starts the liquid moving. The walls prevent it from leaving
the container, so the motion turns into a circulation inside the square.
This is the \emph{lid-driven cavity} problem.

To calculate that motion, we need both the fluid velocity and the pressure
throughout the container. Two physical requirements determine them together.

\textbf{Mass conservation.} Imagine a small fixed region inside the liquid.
For an incompressible liquid, whose density stays constant, the amount entering that region must equal the
amount leaving it. Otherwise liquid would accumulate or disappear there.
This requirement is called the \emph{continuity equation}.

\textbf{Momentum conservation.} The liquid changes its motion because of
forces. Pressure differences push it, and viscosity transfers motion between
neighboring fluid layers. The \emph{momentum equations} relate those effects
to the fluid acceleration.

Neither requirement alone gives a complete solution. A velocity calculated
with an inaccurate pressure can satisfy a momentum calculation while still
giving unequal inflow and outflow in a small region. We therefore need a
procedure that makes velocity and pressure satisfy both requirements.

\textbf{The route through the lesson.}
\begin{enumerate}
\item Write the continuity and momentum equations for this cavity.
\item Divide the domain into small cells and explain the finite-volume method:
      each cell must obey its own conservation balance.
\item Locate the stored pressure and velocity values and derive the equations
      that connect neighboring cells.
\item Explain how an initial pressure estimate is corrected. Then assemble
      these operations into SIMPLE, PISO and PIMPLE.
\item Examine calculated cavity fields, accuracy, convergence and measured cost.
\end{enumerate}
The first calculation uses a Reynolds number
of 100; its definition follows with the governing equations.
